% =========================================================
% Tagged Partitions — Notes (stub, expand before Riemann)
% Chapter: continuity
% Volume III · Analysis
% =========================================================

% ---------------------------------------------------------
% TODO: Expand this section before writing Riemann integration.
% Vocabulary defined here is shared with the Riemann chapter.
% ---------------------------------------------------------

\subsection{Tagged Partitions and Mesh}
\label{sec:tagged-partitions}

% ---------------------------------------------------------
% Definition — Mesh of a Partition
% ---------------------------------------------------------

\begin{tcolorbox}[colback=defbox, colframe=defborder, arc=2pt,
  left=6pt, right=6pt, top=4pt, bottom=4pt,
  title={\small\textbf{Definition (Mesh of a Partition)}},
  fonttitle=\small\bfseries]
\begin{definition}[Mesh of a Partition]
\label{def:mesh-of-partition}
The \textbf{mesh} (or \textbf{norm}) of a partition
$P=\{x_0,\ldots,x_n\}$ of $[a,b]$ is
\[
  \|P\| := \max_{1\leq i\leq n}(x_i - x_{i-1}).
\]
\end{definition}
\end{tcolorbox}


\begin{remark*}[Standard quantified statement]
\[
  \|P\| = \max\{x_i - x_{i-1} : 1\leq i\leq n\}.
\]
\end{remark*}

\begin{remark*}[Interpretation]
The mesh is the width of the widest subinterval. A partition has
small mesh iff every subinterval is narrow. The Riemann integral
is defined by taking limits as $\|P\|\to 0$.
\end{remark*}

% ---------------------------------------------------------
% Definition — Refinement
% ---------------------------------------------------------

\begin{tcolorbox}[colback=defbox, colframe=defborder, arc=2pt,
  left=6pt, right=6pt, top=4pt, bottom=4pt,
  title={\small\textbf{Definition (Refinement)}},
  fonttitle=\small\bfseries]
\begin{definition}[Refinement]
\label{def:refinement-of-partition}
A partition $Q$ is a \textbf{refinement} of $P$ if every point of
$P$ is also a point of $Q$:
\[
  \operatorname{Refinement}(Q,P)
  \;\iff\;
  P \subseteq Q \text{ (as sets of points)}.
\]
\end{definition}
\end{tcolorbox}


\begin{remark*}[Standard quantified statement]
\[
  \operatorname{Refinement}(Q,P)
  \;\iff\;
  \forall x\in P: x\in Q.
\]
\end{remark*}

\begin{remark*}[Interpretation]
Refinement adds partition points without removing any. Every
refinement satisfies $\|Q\|\leq\|P\|$. The Darboux approach to
Riemann integration uses the fact that upper sums decrease and
lower sums increase under refinement.
\end{remark*}

% ---------------------------------------------------------
% Definition — Common Refinement
% ---------------------------------------------------------

\begin{tcolorbox}[colback=defbox, colframe=defborder, arc=2pt,
  left=6pt, right=6pt, top=4pt, bottom=4pt,
  title={\small\textbf{Definition (Common Refinement)}},
  fonttitle=\small\bfseries]
\begin{definition}[Common Refinement]
\label{def:common-refinement}
The \textbf{common refinement} of partitions $P$ and $Q$ of $[a,b]$
is $P\cup Q$ (as sets of points).
\end{definition}
\end{tcolorbox}


\begin{remark*}[Interpretation]
$P\cup Q$ is the coarsest partition that refines both $P$ and $Q$.
Used in Darboux integration proofs to compare upper and lower sums
from two different partitions.
\end{remark*}
